% id: 65-30
% book: 베이직쎈 확률과 통계 · 04 조건부확률
% section: 개념 19 $\mathrm{P}(B)=\mathrm{P}(A\cap B)+\mathrm{P}(\comp{A}\cap B)$를 이용한 확률
% figure: none
% answer: 
% answer_source: 
% variant_level: 0 (원본 전사)
% 이 파일은 build.mjs 가 items.json 에서 만든다. 고칠 때는 items.json 을 고친다.
\smash{\raisebox{1.5mm}{\dmkichul{베이직쎈 확률과 통계 65쪽 30번}}}
2개의 당첨 제비를 포함하여 8개의 제비가 들어 있는 상자에서 임의로 제비를 한 개씩 두 번 뽑을 때, 첫 번째에 당첨 제비를 뽑는 사건을 $A$, 두 번째에 당첨 제비를 뽑는 사건을 $B$라 하자. 다음을 구하시오. \dmcondfill{뽑은 제비는 다시 넣지 않는다.}{}
\dmsubprob{1} 두 번 모두 당첨 제비를 뽑을 확률\dmhint{$\mathrm{P}(A\cap B)=\mathrm{P}(A)\mathrm{P}(\blank[3.2em])$\\ $\hphantom{\mathrm{P}(A\cap B)}=\dfrac{2}{8}\cdot\blank=\blank$}
\dmsubprob{2} 두 번째에만 당첨 제비를 뽑을 확률\dmhint{$\mathrm{P}(\comp{A}\cap B)=\mathrm{P}(\comp{A})\mathrm{P}(\blank[3.2em])$\\ $\hphantom{\mathrm{P}(\comp{A}\cap B)}=\dfrac{6}{8}\cdot\blank=\blank$}
\dmsubprob{3} 두 번째에 당첨 제비를 뽑을 확률\dmhint{두 사건 $A\cap B$와 $\comp{A}\cap B$는 서로 배반사건이므로\\ $\mathrm{P}(B)=\mathrm{P}(A\cap B)+\mathrm{P}(\blank[4em])$\\ $\hphantom{\mathrm{P}(B)}=\blank+\dfrac{3}{14}=\blank$}
